\section{From Workaround to Interface}

\begin{table}[t]
\centering
\small
\begin{tabular}{L{0.20\linewidth} L{0.34\linewidth} L{0.36\linewidth}}
\toprule
Dimension & Source paper (2018) & Upgraded paper (2026 framing) \\
\midrule
Primary question & Can 2D and 1D CNNs compete with 3D CNNs after Hilbert-based mapping? & When should serialization be treated as a first-class interface for irregular 3D and biomolecular learning? \\
Paper type & Empirical methods note & Hybrid interface paper with an executable probe \\
Evidence base & ModelNet10 and K-Ras / H-Ras experiments & Source-paper analysis, modern primary literature, and a deterministic context-budget study \\
Main contribution & A reversible 3D-to-2D / 1D mapping strategy & A taxonomy, executable locality study, design framework, and biomolecular research agenda for serialized structural learning \\
Acceptance story & Better efficiency than volumetric baselines under a 2018 benchmark setup & A timely interface paper that connects early Hilbert mappings to modern sequence-native 3D learning and demonstrates why ordering quality matters under bounded context \\
\bottomrule
\end{tabular}
\caption{Why the new manuscript is not a facelift. The central claim, paper type, evidence base, and contribution all change materially.}
\label{tab:source-vs-upgrade}
\end{table}

\begin{figure}[t]
\centering
\fbox{%
\begin{minipage}{0.92\linewidth}
\footnotesize
\centering
\textbf{Serialization-first view of structural learning}

\vspace{0.5em}
\textbf{Irregular structure} $\rightarrow$ \textbf{locality-aware serialization} $\rightarrow$ \textbf{sequence-native backbone} $\rightarrow$ \textbf{task layer}

\vspace{0.5em}
\begin{tabular}{c}
point cloud / voxel grid / atomistic field \\
Hilbert order / serialized neighbors / multiscale token order \\
transformer / state-space model / masked autoencoder \\
classification / reconstruction / retrieval \\
\end{tabular}

\vspace{0.5em}
\raggedright
\textbf{Design axes:} locality, token budget, channel semantics, reversibility, and multiscale aggregation.
\end{minipage}}
\caption{The upgraded paper's central view. Serialization is treated as a modeling interface that connects irregular 3D structure to efficient sequence-native learning systems, rather than as a one-off preprocessing trick.}
\label{fig:serialization-view}
\end{figure}


The central reframing of this manuscript is straightforward: serialization should be treated as a modeling interface rather than as a preprocessing workaround. This section explains why that reframing is justified by the literature that appeared after the source paper.

The original paper demonstrated that a Hilbert-style traversal could map voxelized 3D structures into lower-dimensional arrays while preserving enough local organization for CNN classification \citep{corcoran2018spatial}. At the time, the practical benefit was obvious. Volumetric models were expensive, low-resolution, and poorly matched to high-channel structural data. A reversible mapping into 2D or 1D offered a way to borrow mature lower-dimensional architectures without discarding all spatial structure.

If the field had simply moved toward stronger native 3D models, the 2018 story might have aged into a historical curiosity. But that is not what happened. Certainly, native point-based models such as Point Transformer raised the baseline for direct geometric learning \citep{zhao2021pointtransformer}. Yet later work also made serialization itself more explicit. Point Transformer V3 replaced exact neighborhood operations with efficient serialized neighbor mappings to improve scale \citep{wu2024ptv3}. PointMamba used space-filling curves for point tokenization in a linear-complexity state-space model \citep{liang2024pointmamba}. Voxel Mamba treated serialization as the basis for group-free sequence modeling in 3D detection \citep{zhang2024voxelmamba}. NoKSR leveraged serialized orderings as a practical approximation layer for neighborhood retrieval in surface reconstruction \citep{li2025noksr}. Across these settings, ordering is no longer incidental. It actively shapes what computation is feasible.

This is why the source paper deserves reinterpretation. Its durable contribution is not that 2D CNNs once matched a particular 3D baseline on a 2018 benchmark. The durable contribution is that it recognized an underappreciated degree of freedom: one may redesign the representation to expose structural information in a form that a different model family can exploit more efficiently. That intuition aged well even though the original empirical framing did not.

Two aspects of the source paper are especially important in hindsight. First, it emphasized locality preservation. This has remained central because serialization is only useful if nearby tokens continue to provide meaningfully local evidence most of the time. Second, the paper emphasized richer channel capacity than a pure volumetric pipeline could cheaply support. That point is even more relevant in biomolecular AI, where geometry alone is rarely enough. Structural tokens may need atom identity, residue type, electrostatic descriptors, occupancy, or learned latent channels attached to them.

The upgraded perspective thus changes the acceptance story. Instead of asking readers to treat the 2018 experiments as newly competitive, it asks them to see the original method as an early instance of a now-general pattern: ordering is a representational contract between structural data and efficient sequence-native computation. In the remainder of the paper, we make that contract more concrete by pairing the reinterpretation with an executable probe of locality and context-budget exposure.
